% frontmatter/czytam/verano.tex -- las instrucciones para el adulto del
% cuaderno de verano de «Czytam z lupą» (ediciones/czytam-verano.tex, ver
% notes/07-ediciones.md): las de frontmatter/lupa/verano.tex, en polaco,
% con las palabras de frontmatter/czytam/como-usar.tex, que habla del año
% entero y es la página a la que sustituye. Detrás, «Jak czyta detektyw»
% (la misma del libro entero) y, en vez del mapa del año, el mapa del
% verano (\mapaEdicion, preamble.tex). Lo que se dice del niño o de la
% niña no tiene género: «dziecko».
\section*{Jak korzystać z tego zeszytu na lato}

To lato z zeszytu \emph{Czytam z lupą}, trzeciego z serii \emph{Uczę się
czytać} po polsku: jego ostatnie 13 tygodni, w osobnym zeszycie, żeby
czytać po trochu każdego dnia wakacji. Jest dla dziecka, które czyta już
płynnie i teraz uczy się czytać uważnie: teksty mają około 110 słów i
długie zdania, z kilkoma częściami połączonymi słowami \textit{że},
\textit{bo}, \textit{chociaż}, \textit{kiedy}, \textit{jeśli} albo
\textit{żeby}, a codziennie jest w nich coś do odkrycia -- nie to, co
wydaje się, że tekst mówi, tylko to, co mówi naprawdę.

Ma 65 stron, po jednej na każdy dzień od poniedziałku do piątku, i każda
wygląda zawsze tak samo:

\begin{itemize}[leftmargin=6mm, itemsep=1mm]
\item \textbf{Data}, \textbf{Czytanie} i \textbf{Uwagi}: dla dorosłego.
\item \textbf{Tekst dnia}, w zielonej ramce, w akapitach. Tytuł ramki
  mówi, jak go czytać.
\item \textbf{Zadanie}, do którego trzeba wrócić do tekstu. Złota
  zasada: \textbf{odpowiedź zawsze jest w tekście} -- w tekście z tego
  dnia albo, jeśli tak mówi polecenie, w tekstach z poprzednich dni tego
  samego tygodnia. Jeśli dziecko jej nie znajduje, nie mów mu jej:
  zapytaj „gdzie to jest napisane?” i pozwól mu przeczytać jeszcze raz.
\end{itemize}

\textbf{Zadania.} \textbf{\lblDibujaConLupa}: narysować dokładnie to, co
opisuje tekst; nie chodzi o to, żeby rysunek był ładny, tylko żeby był
dokładny (ile, jakiego koloru, gdzie), a na koniec sprawdza się go z
tekstem w ręku, odpowiadając na pytania pod spodem. \textbf{\lblMapa}:
iść za wskazówkami po kratkach i zaznaczyć miejsce albo drogę.
\textbf{\lblLogica}: zagadka logiczna; w każdej kratce stawia się ✓ albo
✗, aż zostanie tylko jedna możliwość. \textbf{\lblErrores}: streszczenie
tekstu z błędami, które trzeba znaleźć i poprawić.
\textbf{\lblVerdaderoFalso}: a jeśli fałsz, to jak jest naprawdę.
\textbf{\lblCodigo}: odczytać wiadomość szyfrem, który wyjaśnia tekst.
\textbf{\lblResponde}, \textbf{\lblOrdena}, \textbf{\lblRelaciona},
\textbf{\lblFicha}, \textbf{\lblCompara}, \textbf{\lblAdivina},
\textbf{\lblVinetas} i \textbf{\lblCrea}: pytania „dlaczego?” i „skąd to
wiesz?”, zdarzenia do ułożenia po kolei, informacje do wypisania z
tekstu popularnonaukowego, dwie rzeczy do porównania, zagadki, historia
w trzech rysunkach i trochę własnego pisania.

\textbf{Co tydzień zagadka.} Strony opowiadają lato \textbf{Lucíi}, jej
brata \textbf{Daniego} i ich psa \textbf{Toby'ego} na wsi u babci Rosy.
Lucía ma klub detektywów, \textbf{Klub Lupy}, ze swoimi przyjaciółmi,
Sofíą i Hugiem; ale oni spędzają lato gdzie indziej, więc na wsi klub to
Lucía, Dani i Toby -- i Martín, kolega, którego Dani ma na wsi. Prawie
co tydzień jest tajemnica: wskazówki pojawiają się od poniedziałku do
czwartku, a w piątek trzeba ją rozwiązać -- \textbf{\lblCaso}: zakreślić
odpowiedź i wypisać wskazówki, które jej dowodzą. Rozwiązanie
potwierdza się w następny poniedziałek, a jest też w
\textbf{odpowiedziach} na końcu zeszytu. Lepiej tam nie zaglądać, zanim
dziecko samo spróbuje. A wielka tajemnica lata to mapa skarbu, którą
babcia Rosa narysowała, kiedy miała osiem lat.

\textbf{Tempo.} Dziesięć albo piętnaście minut dziennie. Jeśli tekst
wydaje się dziecku za długi, przeczytajcie go najpierw razem na głos, po
jednym akapicie, a potem niech przeczyta go jeszcze raz samo.
Najważniejsze, żeby to wciąż była zabawa: detektywi się nie spieszą,
detektywi są uważni. Na następnej stronie jest to, jak czyta detektyw, a
potem lato tydzień po tygodniu, ze słońcem do pokolorowania za każdy
dzień czytania; na końcu -- dyplom.
\newpage
